\chapter{The pilot session, in full}
\label{app:pilot}

\noindent\emph{\Cref{sec:pilot-results} reports the headline pilot numbers.
This appendix is the full record: coverage, every session, and the
supplementary measures.}

Twenty-three sessions were recorded, sixteen on 2026-09-03/04 for P1--P4 and
seven on 2026-09-02 for P5 (\repo{recordings/sessions/}). Duration, the
robot's end-effector path and lag come from the full-resolution position trace
(20 Hz); clutch and emergency-stop counts come from each session's event
log. The session folders were named at the console, and the names mixed
the two tasks up: sessions named "object tracking" were the pick-and-place
task, and "position matching" or "position reaching" was the same
target-reaching task under another name. The operator who ran the sessions
confirmed the mapping, and every table and figure here uses the two real
task names (\repo{figures/relabel\_pilot\_tasks.py}). P5's sessions predate the per-participant folder-naming convention
(generic timestamp names, no task label) and are identified as direct or
shared from the \texttt{shared\_autonomy} flag each session's own
\texttt{events.jsonl} logs at start, which -- unlike the folder names relied
on for P1--P4 -- is populated correctly for every one of P5's sessions.
Participant identity is P1--P5 throughout, including in every figure pixel,
consistent with \repo{write\_manifest()} refusing to record a name
(\cref{sec:ethics}).

\begin{table}[htbp]
  \centering
  \small
  \caption[Pilot coverage]{Which participant performed which task under which
  condition. Coverage is uneven because a pilot session is limited by time.
  P5's task was not recorded; all seven of P5's sessions are one
  undifferentiated block.}
  \label{tab:pilot-coverage}
  \begin{tabular}{@{}lccc@{}}
    \toprule
    Participant & Pick and place & Target reaching & (unlabelled) \\
    \midrule
    P1 & both (direct $\times2$) & --- & --- \\
    P2 & both & both (shared $\times2$) & --- \\
    P3 & direct only & both (shared $\times2$) & --- \\
    P4 & both & both & --- \\
    P5 & --- & --- & both (direct $\times5$, shared $\times2$) \\
    \bottomrule
  \end{tabular}
\end{table}

\begin{table}[htbp]
  \centering
  \scriptsize
  \caption[Every pilot session]{Every recorded session. Duration includes
  setup and hesitation. Active share is the fraction of the session spent
  with the grip engaged. P1's two direct-control sessions are a short
  warm-up followed by the recorded attempt, both retained as separate
  points.}
  \label{tab:pilot-sessions}
  \begin{tabular}{@{}llrrrrr@{}}
    \toprule
    P & Task & Cond. & Duration (s) & Travel (m) & Re-grips & Active \\
    \midrule
    P1 & Pick and place   & Direct & 147.4 & 0.31 & 19 & \SI{50}{\percent} \\
    P1 & Pick and place   & Direct & 392.8 & 1.85 & 53 & \SI{53}{\percent} \\
    P1 & Pick and place   & Shared & 114.9 & 1.50 & 12 & \SI{61}{\percent} \\
    P2 & Pick and place   & Direct & 38.7  & 0.90 & 2  & \SI{61}{\percent} \\
    P2 & Pick and place   & Shared & 169.6 & 1.57 & 25 & \SI{76}{\percent} \\
    P2 & Target reaching  & Direct & 227.9 & 2.77 & 25 & \SI{52}{\percent} \\
    P2 & Target reaching  & Shared & 78.6  & 0.70 & 11 & \SI{65}{\percent} \\
    P2 & Target reaching  & Shared & 357.9 & 5.39 & 41 & \SI{73}{\percent} \\
    P3 & Pick and place   & Direct & 617.4 & 7.09 & 52 & \SI{83}{\percent} \\
    P3 & Target reaching  & Shared & 458.1 & 3.18 & 51 & \SI{74}{\percent} \\
    P3 & Target reaching  & Direct & 315.9 & 3.72 & 35 & \SI{67}{\percent} \\
    P3 & Target reaching  & Shared & 196.8 & 4.08 & 14 & \SI{85}{\percent} \\
    P4 & Pick and place   & Direct & 195.6 & 2.07 & 37 & \SI{68}{\percent} \\
    P4 & Pick and place   & Shared & 186.9 & 1.98 & 25 & \SI{42}{\percent} \\
    P4 & Target reaching  & Direct & 171.9 & 1.65 & 12 & \SI{30}{\percent} \\
    P4 & Target reaching  & Shared & 111.7 & 1.98 & 12 & \SI{40}{\percent} \\
    P5 & (unlabelled) & Direct & 262.5 & 1.99 & 20 & \SI{37}{\percent} \\
    P5 & (unlabelled) & Shared & 250.0 & 1.58 & 23 & \SI{39}{\percent} \\
    P5 & (unlabelled) & Direct & 359.5 & 5.19 & 43 & \SI{62}{\percent} \\
    P5 & (unlabelled) & Shared & 265.8 & 5.89 & 19 & \SI{53}{\percent} \\
    P5 & (unlabelled) & Direct & 207.8 & 1.62 & 15 & \SI{43}{\percent} \\
    P5 & (unlabelled) & Direct & 93.0  & 0.16 & 11 & \SI{56}{\percent} \\
    P5 & (unlabelled) & Direct & 99.3  & 0.59 & 11 & \SI{57}{\percent} \\
    \bottomrule
  \end{tabular}
\end{table}

\begin{figure}[htbp]
  \centering
  \begin{subfigure}[b]{0.32\linewidth}
    \centering\includegraphics[width=\linewidth]{pilot/pilot_distance.pdf}
    \caption{robot end-effector path}
  \end{subfigure}\hfill
  \begin{subfigure}[b]{0.32\linewidth}
    \centering\includegraphics[width=\linewidth]{pilot/pilot_regrips.pdf}
    \caption{re-grips}
  \end{subfigure}\hfill
  \begin{subfigure}[b]{0.32\linewidth}
    \centering\includegraphics[width=\linewidth]{pilot/pilot_active.pdf}
    \caption{active share}
  \end{subfigure}\\[6pt]
  \begin{subfigure}[b]{0.32\linewidth}
    \centering\includegraphics[width=\linewidth]{pilot/pilot_lag.pdf}
    \caption{control lag}
  \end{subfigure}\hfill
  \begin{subfigure}[b]{0.60\linewidth}
    \centering\includegraphics[width=\linewidth]{pilot/pilot_paired_regrips.pdf}
    \caption{re-grips, paired within operator and task}
  \end{subfigure}
  \caption[Pilot: supplementary measures]{The supplementary pilot measures,
  by task and condition. Lag is essentially unchanged between conditions;
  the robot's end-effector path (a) is longer under shared autonomy in two
  of three task groups, consistent with the paired result in
  \cref{sec:pilot-limits} -- and it is the \emph{robot's} path: the
  operators' own controller paths differ by \SI{2}{\percent} between
  conditions (\SI{29.9}{\metre} direct, \SI{30.5}{\metre} shared). Both
  are drawn from the position trace rather than a stopwatch, which is why
  they are read as more reliable than session duration in a pilot where
  duration is confounded with setup time.}
  \label{fig:pilot-supplementary}
\end{figure}

\begin{figure}[htbp]
  \centering
  \begin{subfigure}[b]{0.30\linewidth}
    \centering\includegraphics[width=\linewidth]{pilot/pilot_traj_direct.pdf}
    \caption{direct control}
  \end{subfigure}\hfill
  \begin{subfigure}[b]{0.30\linewidth}
    \centering\includegraphics[width=\linewidth]{pilot/pilot_traj_shared.pdf}
    \caption{shared autonomy}
  \end{subfigure}\hfill
  \begin{subfigure}[b]{0.36\linewidth}
    \centering\includegraphics[width=\linewidth]{pilot/pilot_engagement.pdf}
    \caption{active-control span, same session}
  \end{subfigure}
  \caption[Pilot: recorded path and engagement, one participant]{The actual
  recorded end-effector path (\repo{ee\_left\_x/y/z} in \repo{trail.csv}) for
  one participant on the same task under both conditions, recorded back to
  back, with the corresponding span of active control alongside. Neither
  path is smooth; a shorter, straighter line would indicate a more direct
  approach to the target. Gaps in the engagement panel are the operator
  releasing the grip to reposition their hand before re-engaging.}
  \label{fig:pilot-trajectory}
\end{figure}

\section{What the gripper record says about pick and place}
\label{sec:gripper-record}

The \SI{20}{\hertz} trail never logged a gripper command (its
\repo{grip\_cmd\_*} columns are empty in every session), so whether a pick
happened was read from two other records (\cref{tab:gripper-use}): the
controller's index trigger, which is what the live gripper node closes the
hand from, passed through that node's own latch; and the physical gripper's
measured finger position from the bag. The two agree. P1 closed the gripper
eight times and held it closed for \SI{18}{\second} in total, P2 four times
for \SI{5.8}{\second}, and P5 four times for \SI{133}{\second} -- carrying
the object; P3 closed it once, to \num{0.90} of full travel, for
\SI{0.7}{\second}; P4 never closed it in any session. That is the record
behind the statement in \cref{sec:pilot-results}: three pick-and-places
completed, one lift, one reach. P5's task labels were never recorded, and
the two long holds identify which of P5's sessions were the pick-and-place
attempts.

\begin{table}[htbp]
  \centering
  \scriptsize
  \caption[Gripper use in every VR session]{Gripper use in every VR
  session, from two independent records. \emph{Commanded}: closures of the
  controller's index trigger in the \SI{20}{\hertz} trail, counted through
  the live gripper node's own latch (close above 0.30, release below 0.10),
  and the share of the session the command was held. \emph{Measured}: the
  physical gripper's own finger position from the bag
  (\repo{/real/gripper\_<hand>}, 0 open, 1 closed): closures, share of the
  session closed, and the furthest the fingers closed. The active hand only.}
  \label{tab:gripper-use}
  \begin{tabular}{@{}lllrrrrr@{}}
    \toprule
    & & & \multicolumn{2}{c}{Commanded} & \multicolumn{3}{c}{Measured} \\
    \cmidrule(lr){4-5} \cmidrule(lr){6-8}
    P & Task & Cond. & Closures & Held & Closures & Closed & Max \\
    \midrule
    P1 & Pick and place & Direct & 2 & \SI{0.4}{\percent} & 2 & \SI{1.3}{\percent} & 1.00 \\
    P1 & Pick and place & Direct & 2 & \SI{0.2}{\percent} & 2 & \SI{0.6}{\percent} & 0.99 \\
    P1 & Pick and place & Shared & 6 & \SI{37.1}{\percent} & 4 & \SI{12.0}{\percent} & 0.99 \\
    P2 & Pick and place & Direct & 1 & \SI{2.7}{\percent} & 0 & \SI{0.0}{\percent} & 0.01 \\
    P2 & Pick and place & Shared & 3 & \SI{17.8}{\percent} & 2 & \SI{2.0}{\percent} & 0.99 \\
    P2 & Target reaching & Direct & 1 & \SI{0.0}{\percent} & 1 & \SI{0.5}{\percent} & 0.93 \\
    P2 & Target reaching & Shared & 1 & \SI{1.1}{\percent} & 1 & \SI{1.8}{\percent} & 1.00 \\
    P2 & Target reaching & Shared & 0 & \SI{0.0}{\percent} & 0 & \SI{0.0}{\percent} & 0.01 \\
    P3 & Target reaching & Direct & 0 & \SI{0.0}{\percent} & 0 & \SI{0.0}{\percent} & 0.01 \\
    P3 & Pick and place & Shared & 2 & \SI{18.2}{\percent} & 2 & \SI{18.5}{\percent} & 0.94 \\
    P3 & Target reaching & Direct & 0 & \SI{0.0}{\percent} & 0 & \SI{0.0}{\percent} & 0.01 \\
    P3 & Target reaching & Shared & 1 & \SI{0.1}{\percent} & 1 & \SI{1.5}{\percent} & 1.00 \\
    P3 & Pick and place & Direct & 1 & \SI{41.1}{\percent} & 1 & \SI{39.9}{\percent} & 0.95 \\
    P3 & Target reaching & Direct & 0 & \SI{0.0}{\percent} & 0 & \SI{0.0}{\percent} & 0.01 \\
    P3 & Target reaching & Direct & 0 & \SI{0.0}{\percent} & 0 & \SI{0.0}{\percent} & 0.01 \\
    P4 & Pick and place & Direct & 2 & \SI{0.1}{\percent} & 1 & \SI{0.1}{\percent} & 0.90 \\
    P4 & Target reaching & Shared & 0 & \SI{0.0}{\percent} & 0 & \SI{0.0}{\percent} & 0.01 \\
    P4 & Target reaching & Direct & 0 & \SI{0.0}{\percent} & 0 & \SI{0.0}{\percent} & 0.01 \\
    P4 & Target reaching & Shared & 0 & \SI{0.0}{\percent} & 0 & \SI{0.0}{\percent} & 0.01 \\
    P5 & Pick and place & Direct & 0 & \SI{0.0}{\percent} & 0 & \SI{0.0}{\percent} & 0.00 \\
    P5 & Pick and place & Shared & 0 & \SI{0.0}{\percent} & 0 & \SI{0.0}{\percent} & 0.00 \\
    P5 & Target reaching & Direct & 0 & \SI{0.0}{\percent} & 0 & \SI{0.0}{\percent} & 0.00 \\
    P5 & Target reaching & Shared & 0 & \SI{0.0}{\percent} & 0 & \SI{0.0}{\percent} & 0.00 \\
    \bottomrule
  \end{tabular}
\end{table}


\section{Sim-to-real tracking during the pilot}

Every pilot session polled the physical Kinova arms' own joint encoders
alongside the simulated joint state, which makes a direct check possible:
does the real arm actually follow the simulated one it was commanded
through. \Cref{fig:track} is the best-tracking session across all twenty-three
(P4, pick and place, direct control, chosen by lowest sim-to-real joint
RMS): the real arm's joint 4 angle follows the simulated one to
\SI{1.06}{\degree} RMS over \num{3859} samples, with the largest departures
during the fastest motion and near-zero error while the arm holds still.
This is the same terminal-offset mechanism \cref{sec:sim-to-real} measures
from replayed commands, seen here live during an operator session rather
than in an offline replay -- \textbf{it is evidence the sim-to-real bridge
was live and tracking throughout the pilot}, not a claim that any pilot
result was validated against a loaded arm, which requires the arm to be
carrying and moving an object rather than only mirroring a joint angle.

\begin{figure}[htbp]
  \centering
  \begin{subfigure}[b]{0.31\linewidth}
    \centering\includegraphics[width=\linewidth]{pilot/track_joint.pdf}
    \caption{joint 4, sim vs.\ real}
    \label{fig:track-a}
  \end{subfigure}\hfill
  \begin{subfigure}[b]{0.31\linewidth}
    \centering\includegraphics[width=\linewidth]{pilot/track_joint_error.pdf}
    \caption{sim-real tracking error}
    \label{fig:track-b}
  \end{subfigure}\hfill
  \begin{subfigure}[b]{0.34\linewidth}
    \centering\includegraphics[width=\linewidth]{pilot/track_cartesian.pdf}
    \caption{operator's hand vs.\ robot hand}
    \label{fig:track-c}
  \end{subfigure}
  \caption[Sim-to-real tracking, best pilot session]{The best-tracking pilot
  session by sim-to-real joint RMS. (\subref{fig:track-a}) One joint,
  simulated command against the real arm's own encoder, over the session.
  (\subref{fig:track-b}) The resulting tracking error, root-mean-square over
  all seven joints; the dotted line is the session RMS,
  \SI{1.06}{\degree}. (\subref{fig:track-c}) The operator's raw VR-hand path
  against the simulated robot hand's achieved path, each centred on its own
  mean -- the two live in different coordinate frames and at different
  scales by design, so an absolute overlay would not be meaningful; centring
  each on itself is the honest way to compare their shapes.}
  \label{fig:track}
\end{figure}

\section{The master-arm cohort}

Three further operators (M1--M3, disjoint from P1--P4) attempted the same
protocol from the instrumented mannequin master rather than VR controllers,
across sixteen further recording attempts on 2026-09-02 and 2026-09-04.
\textbf{Seven of those sixteen are confirmed empty}: \texttt{ros2 bag info}
reports zero messages, or no bag metadata at all, for every attempt at the
multimeter task and both of M1's target-reaching attempts -- the same
recording-pipeline failure this project's own session index for that day
logged independently (eight of twenty-six recordings clean, ten with no data
at all). The nine that did record are real sessions,
\SIrange{72}{304}{\second} each, tens of thousands of messages.

\begin{table}[htbp]
  \centering
  \small
  \caption[Master-arm sessions]{The nine recorded master-arm sessions.
  Joint RMS is the sim-to-real tracking error, computed the same way as
  \cref{fig:track}, over the active arm's seven joints; see
  \cref{fig:master-tracking}.}
  \label{tab:master-sessions}
  \begin{tabular}{@{}llrrr@{}}
    \toprule
    M & Task, condition & Duration (s) & Master path (m) & Joint RMS (deg) \\
    \midrule
    M1 & Pick and place, direct   & 154.1 & 21.44 & 21.2 \\
    M1 & Pick and place, direct   & 145.8 & 14.11 & 15.9 \\
    M1 & (unspecified), shared     & 209.8 & 22.02 & 16.3 \\
    M2 & Pick and place, direct   & 110.9 & 10.85 & 2.2 \\
    M2 & Pick and place, shared   & 71.7  & 7.57  & 1.9 \\
    M2 & Target reaching, direct   & 303.8 & 20.80 & 1.8 \\
    M2 & Target reaching, shared   & 83.6  & 4.19  & 0.4 \\
    M2 & Target reaching, shared   & 206.4 & 11.60 & 3.0 \\
    M3 & Target reaching, direct   & 254.1 & 11.48 & 1.8 \\
    \bottomrule
  \end{tabular}
\end{table}

\begin{figure}[htbp]
  \centering
  \includegraphics[width=0.72\linewidth]{pilot/master_tracking.pdf}
  \caption[Master-arm sim-to-real tracking, all nine sessions]{Sim-to-real
  joint tracking error for every recorded master-arm session, sorted low to
  high. M2 and M3 track within \SI{3}{\degree} throughout; \textbf{every one
  of M1's three sessions tracks an order of magnitude worse}
  (\SIrange{15.9}{21.2}{\degree}).}
  \label{fig:master-tracking}
\end{figure}

M1's tracking error is not noise: it is consistent across all three of that
operator's sessions and separated from M2 and M3 by a factor of five to
fifty. The mechanism this project's own channel-health measurement predicts
(\cref{ch:results}) is a session that leans on the incoherent left-arm
channels rather than the coherent right-arm ones -- all three of M1's
sessions used the \emph{right} arm as the active hand (\cref{tab:master-sessions}),
so a simple left/right split does not explain it on its own, and which
specific channels M1's sessions exercised has not been isolated. It is
reported here as a finding this pilot surfaced, not a diagnosis: the
master-arm modality can apparently produce sim-to-real tracking an order of
magnitude worse than either VR teleoperation or another operator on the same
hardware, on the same day, and the pooled VR-only comparison of
\cref{sec:pilot-results} would have missed it entirely.

\paragraph{Limits.} Five participants with uneven task coverage is a
pilot-scale sample: no significance test is reported, and none of the
differences above should be read as a proven effect. The pooled figures mix
the two tasks together; the paired panels are the part of this pilot closest
to the paired design of \cref{app:study}. Session duration runs from the
start of the recording to the end, including any explanation or hesitation,
so it is a session-length measure rather than a pure motor-execution time.
Every session's emergency stop was engaged once, near the end, to stop the
recording -- a procedural marker rather than a safety incident -- and is
excluded throughout. This pilot has no separate wearer: every measure here
is what the operator experienced, and the wearer-side hypotheses of
\cref{app:study} (H3, H4) remain untested.
